\documentclass[11pt,a4paper]{article}
\usepackage[T1]{fontenc}
\usepackage[utf8]{inputenc}
\usepackage[english]{babel}
\usepackage{amsmath,amssymb}
\usepackage[margin=2cm]{geometry}
\setlength{\parindent}{0pt}
\setlength{\emergencystretch}{3em}
\allowdisplaybreaks

\title{About \texttt{book\_structures\_beyond.tex}}
\author{}
\date{}
\begin{document}
\maketitle
This note explains what the test data in \texttt{book\_structures\_beyond.tex} is used for. That file is input for the engine in \texttt{main.c}; it is read by the test \texttt{test\_algebraic\_structures.c} in this folder.

\begin{enumerate}
\item \textbf{Purpose.} The file holds 65 statements from the same textbook exercises as \texttt{book\_structures.tex} that the engine does not decide yet, or whose notation it cannot read.

\item \textbf{What it exercises.} Honesty about the limits: the engine may answer that it cannot decide, but must never call one of these statements false. When one becomes provable it should move to \texttt{book\_structures.tex}.

\item \textbf{Variety.} An exponent that is a letter; greatest common divisors, least common multiples and primes; the closed part generated by a set; the composition of maps as an operation and maps to be built; residue classes; quotients and infinite sets; a few congruences modulo a letter whose proof needs more than a combination of the hypotheses.

\item \textbf{Expected result.} The line \texttt{PASS: none of the 65 true statements \ldots is called false}.

\item \textbf{How to run it.} From the top folder of the repository:
\begin{verbatim}
cc -std=c11 -O2 -Wall -Wextra -o /tmp/test_algebraic_structures \
    tests/algebraic_structures/test_algebraic_structures.c \
    && /tmp/test_algebraic_structures
\end{verbatim}
\end{enumerate}
\end{document}
